\documentclass[11pt]{article}
\usepackage[a4paper,margin=1in]{geometry}
\usepackage{amsmath,amssymb}
\usepackage{booktabs}
\usepackage[hidelinks]{hyperref}
\usepackage{xcolor}

\setlength{\parskip}{0.4em}
\setlength{\parindent}{0pt}

\providecommand{\rupee}{\textrm{Rs.}}

\newcommand{\intuit}{\smallskip\noindent\textbf{The plain idea.}\ }
\newcommand{\example}{\smallskip\noindent\textbf{A worked number example.}\ }
\newcommand{\formula}{\smallskip\noindent\textbf{The equation.}\ }
\newcommand{\vars}{\smallskip\noindent\textbf{Every term, and what kind of thing it is.}\ }
\newcommand{\applied}{\smallskip\noindent\textbf{How I used it in my project.}\ }
\newcommand{\tells}{\smallskip\noindent\textbf{What it tells you.}\ }

\title{\textbf{The Theory and the Maths}\\[4pt]
\large A from-scratch, plain-language guide to every method behind my evaluation of\\
India's 2021 agricultural-derivatives suspension}
\author{}
\date{}

\begin{document}
\maketitle

\noindent\textbf{Status note, updated 6 July 2026.} The worked results in this companion are
stated on the current \emph{food-donor} rerun: aggregate DiD about $-9.8\%$, not statistically
significant under few-cluster inference (CR1 $p \approx 0.145$; wild bootstrap $p \approx 0.153$),
SDID $-19.1\%$ ($z=-1.88$), with the strongest negative signal concentrated in chana. The older
clean-core $-20.7\%$ run appears only as the explicitly labelled ``before'' stage of the
donor-sensitivity arc (it is part of the story, not the result).

\begin{abstract}
\noindent This is the maths companion to my internship project, which asks whether India's
December 2021 suspension of exchange-traded agricultural derivatives actually calmed food
prices. I teach every quantitative method I used, from the ground up, in plain words. Each
section follows the same shape: the plain idea, a worked example with real numbers, the
equation with \emph{every} symbol named and classified (which terms are the data, which are
the things I estimate, which is the ``dependent'' variable being explained and which are the
``independent'' variables doing the explaining), and finally exactly how I used it in the
project and what it told me. You need only a little prior statistics. By the end you should be
able to read every number in my results and know precisely how it was built.
\end{abstract}

\tableofcontents
\bigskip

% =====================================================================
\section{The setting, in one paragraph}

On 20 December 2021 India's market regulator --- SEBI, the Securities and Exchange Board of
India --- switched off exchange-traded \emph{derivatives} (financial contracts whose value is
tied to an underlying commodity; here \emph{futures} --- agreements to buy or sell a crop at a
price fixed today for delivery on a future date) in seven food commodities, hoping to calm food
inflation. My job was to measure whether the cash (\emph{spot}) market for those crops became
calmer or more turbulent afterwards. The crops that were banned (wheat, chana/chickpea, mustard,
soybean, paddy, moong, crude palm oil) are my \emph{treated} group; a set of crops that kept
trading (guar seed, castor, jeera/cumin, turmeric, cotton) are my \emph{control} group. Every
method below exists to answer one question: did switching off futures change how much spot
prices \emph{jump around}, and can I trust the answer?

\bigskip
\noindent\textbf{Two words you will need throughout.} A \emph{dependent} variable is the thing
I am trying to explain (the outcome, what I put on the left of an equation). An
\emph{independent} variable is something I use to explain it (a driver, what I put on the right).
A \emph{parameter} is a number I do not observe directly but \emph{estimate} from the data (the
strength of a relationship). Keeping these three straight --- outcome, driver, estimated number
--- is most of what reading these equations comes down to.

% =====================================================================
\section{Returns and realized volatility}

This is the raw material for everything else. Before I can ask whether prices got more
``jumpy,'' I need a precise number for jumpiness.

\subsection{Daily log returns}

\intuit A \emph{return} is the percentage change in a price from one day to the next. Instead of
the plain percentage I use the \emph{log return} --- the natural logarithm of the price ratio.
Two reasons. First, log returns \emph{add up cleanly over time}: the log return over a week is
just the sum of the seven daily log returns, which plain percentages do not give you. Second,
they are \emph{symmetric}: a price that rises by a factor and then falls by the same factor
returns exactly to where it started, because $+x$ and $-x$ in logs cancel, whereas a $+10\%$
move followed by a $-10\%$ move leaves you slightly poorer in plain percentages.

\example Suppose a mandi (market) price runs over four days:
\[
P_0 = 100,\quad P_1 = 110,\quad P_2 = 99,\quad P_3 = 104.
\]
The daily log returns are
\[
r_1 = \ln\!\Big(\tfrac{110}{100}\Big) = \ln(1.10) \approx +0.0953,\quad
r_2 = \ln\!\Big(\tfrac{99}{110}\Big) = \ln(0.90) \approx -0.1054,
\]
\[
r_3 = \ln\!\Big(\tfrac{104}{99}\Big) = \ln(1.0505) \approx +0.0493.
\]
So roughly $+9.5\%$, then $-10.5\%$, then $+4.9\%$. Notice the day-1 up-move and the day-2
down-move are almost equal and opposite in logs even though the rupee amounts ($+10$ then $-11$)
are not --- that symmetry is the point. And the total move from day~0 to day~3 is
$r_1+r_2+r_3 \approx +0.0392 = \ln(104/100)$: the daily logs simply sum to the whole-period log
return. They add up.

\formula For a price $P_t$ on day $t$, the daily log return is
\[
r_t \;=\; \ln\!\left(\frac{P_t}{P_{t-1}}\right) \;=\; \ln P_t - \ln P_{t-1}.
\]

\vars
\begin{itemize}
\item $P_t$ --- the price today. \emph{Observed data} (an input). Here, a daily mandi price.
\item $P_{t-1}$ --- the price yesterday. Also observed data.
\item $r_t$ --- the log return. A \emph{derived} number (computed, not estimated): the output of
this step. It is unit-free (a pure percentage-like number), so two commodities priced in very
different rupee amounts become directly comparable.
\item $\ln$ --- the natural logarithm. A fixed mathematical function, not a variable.
\end{itemize}
There is no ``dependent vs independent'' split here yet --- this is just a transformation of the
raw price into a clean building block. That split appears once I run a regression (Section~3).

\tells $r_t$ is a clean, units-free measure of ``how big was today's price move, in percent.'' A
string of these returns is what I feed into the volatility calculation next.

\subsection{Realized volatility}

\intuit \emph{Volatility} (how much a price jumps around day to day) is the \emph{spread} of
those daily returns. If the returns are all tiny and similar, the price is calm; if they swing
between large positives and large negatives, it is volatile. The standard textbook measure of
spread is the \emph{standard deviation} --- the typical distance of the returns from their own
average. ``Realized'' volatility just means I compute it directly from the returns that actually
happened, assuming no model. I do it over a \emph{rolling window} (e.g.\ the last 30 trading
days), so the window slides forward each day and I get a fresh reading for every day. Finally I
\emph{annualise} by multiplying by $\sqrt{252}$ (about 252 trading days in a year) to put the
number on the familiar ``percent per year'' scale.

\example Take three daily returns: $r = \{+0.02,\ -0.01,\ +0.03\}$ (i.e.\ $+2\%, -1\%, +3\%$).
Their average is $\bar r = (0.02 - 0.01 + 0.03)/3 = 0.0133$. The squared deviations from that
average are
\[
(0.02-0.0133)^2 = 0.0000444,\quad
(-0.01-0.0133)^2 = 0.000544,\quad
(0.03-0.0133)^2 = 0.000278.
\]
Summing and dividing by $w-1 = 2$ gives a variance of $0.000433$, so the (daily) standard
deviation is $\sqrt{0.000433} \approx 0.0208$, about $2.1\%$ per day. Annualised, that is
$0.0208 \times \sqrt{252} \approx 0.330$, i.e.\ roughly $33\%$ per year. That single number,
$33\%$, is the realized volatility for this little stretch.

\formula Over a rolling window of $w$ trading days ending at day $t$, realized volatility is
\[
\sigma_t \;=\; \sqrt{\frac{1}{w-1}\sum_{i=t-w+1}^{t}\big(r_i - \bar r\big)^2}\;\times\;\sqrt{252}.
\]

\vars
\begin{itemize}
\item $\sigma_t$ --- realized volatility at day $t$. \emph{Derived output}, in percent per year.
This is the number I care about; later it becomes the \emph{dependent} variable in my regressions.
\item $r_i$ --- the daily log returns inside the window. Observed (well, derived from prices) inputs.
\item $\bar r$ --- their average over the window. A derived helper number.
\item $w$ --- the window length: a \emph{choice} I make, not estimated. I use $w=30$, with $w=60$
as a robustness check (a longer window = smoother, slower-moving volatility).
\item $w-1$ --- the divisor (instead of $w$). A small textbook correction so the variance estimate
is not biased downward; ``one degree of freedom used up'' by estimating $\bar r$.
\item $\sqrt{252}$ --- the annualisation factor. A fixed constant, converting daily to yearly scale.
\end{itemize}

\applied I computed $\sigma_t$ this way for every commodity, for every district, for every day,
from the cleaned spot prices --- then averaged to the month to get one volatility number per
commodity-district-month. That monthly $\ln\sigma_{it}$ is the \emph{outcome} I carry into the
difference-in-differences and event-study regressions below. I deliberately kept it
\emph{model-free} (a plain standard deviation, no distribution assumed) because the whole project
hinges on a volatility figure, and I did not want a hidden modelling assumption quietly doing the
work.

\smallskip\noindent\textbf{One honest caveat.} Consecutive 30-day windows share 29 of their 30
returns, so two neighbouring days' readings are almost the same number by construction. The series
is therefore mechanically smooth and strongly correlated from day to day. That \emph{overlap} can
make trends look more real than they are, so I do not over-read the day-by-day path; the 60-day
window and a non-overlapping monthly version are checks against it.

% =====================================================================
\section{Difference-in-Differences (DiD)}

This is the core of the volatility analysis: how I try to isolate the ban's effect from
everything else happening in 2022.

\subsection{The idea: a double difference}

\intuit The naive approach --- compare banned-commodity volatility before the ban to after --- is
hopeless, because 2022 was chaos for \emph{every} commodity: the Ukraine war, an energy spike, and
global food inflation all hit at once. If banned-crop volatility rose, was it the ban or the war?
You cannot tell from a single before/after. The fix is a \emph{control group} of crops that kept
trading futures, and a \emph{double difference}. First difference: the change in volatility (after
minus before) for the banned crops. Second difference: the same change for the control crops. The
DiD estimate is the difference between those two changes. The logic: the controls lived through
exactly the same 2022, so whatever the war did shows up in \emph{both} groups and \emph{cancels}
when I subtract. What is left --- the banned crops' \emph{extra} change beyond the controls --- is
what I attribute to the ban.

\subsection{The $2\times 2$ table}

\example Suppose the (annualised) volatilities are:

\begin{center}
\begin{tabular}{lccc}
\toprule
 & \textbf{Before ban} & \textbf{After ban} & \textbf{Change} \\
\midrule
\textbf{Treated} (banned crops) & 40\% & 50\% & $+10$ \\
\textbf{Control} (still trading) & 30\% & 36\% & $+6$ \\
\midrule
 & & \textbf{DiD} & $+10 - (+6) = +4$ \\
\bottomrule
\end{tabular}
\end{center}

Both groups got more volatile (the war hit everyone), but the treated group rose by $10$ points
and the controls by only $6$. The double difference is $10 - 6 = +4$ percentage points --- the
slice of the rise I assign to the ban, after differencing out the common shock that pushed both
groups up. If instead the controls had \emph{also} risen by $10$, the DiD would be $0$: the ban
added nothing beyond what hit everyone.

\subsection{The panel regression}

\intuit With many commodities, many regions, and many months, I do not literally fill in a
$2\times2$ table; I run a regression that does the same double difference automatically, using
\emph{fixed effects} (dummy variables that soak up whole categories of difference).

\formula
\[
\ln \sigma_{it} \;=\; \alpha_i \;+\; \lambda_t \;+\; \beta\,\big(\mathrm{Treat}_i \times \mathrm{Post}_t\big) \;+\; \varepsilon_{it}.
\]

\vars This is the first equation with a true dependent/independent split, so I will be explicit.
\begin{itemize}
\item $\ln \sigma_{it}$ --- log realized volatility for unit $i$ (a commodity, or a
commodity$\times$district) in month $t$. This is the \textbf{dependent variable} --- the outcome,
the thing I am explaining, on the left-hand side. (I take the log so that the coefficient reads as
a percentage, and so that proportional changes, not rupee changes, are what matter.)
\item $\mathrm{Treat}_i$ --- an \textbf{independent variable}; an \emph{indicator} (dummy) equal to
$1$ for banned crops and $0$ for controls. Property: it is \emph{time-invariant} --- it never
changes over $t$ (a crop is banned or not, always).
\item $\mathrm{Post}_t$ --- an \textbf{independent variable}; an indicator equal to $1$ for months
after 20 December 2021 and $0$ before. Property: it is \emph{common to all units} in a given month.
\item $\mathrm{Treat}_i \times \mathrm{Post}_t$ --- the \textbf{interaction}: $1$ only for a banned
crop in a post-ban month. This is the single independent variable I actually care about.
\item $\beta$ --- a \textbf{parameter} (estimated, not observed): the coefficient on the
interaction. \emph{This is the DiD estimate} --- the one number the whole regression exists to
produce. Because the outcome is in logs, $\beta$ reads approximately as a percent \emph{when it is
small}: $\beta=0.09$ means about $+9\%$. When $\beta$ is not small I convert exactly with
$e^{\beta}-1$ (so my current estimate $\beta=-0.103$ is a $-9.8\%$ effect; and at the earlier
clean-core stage $\beta=-0.232$ was a $-20.7\%$ effect, \emph{not} $-23.2\%$).
\item $\alpha_i$ --- \textbf{unit fixed effects}: one estimated number per unit, absorbing
\emph{permanent} differences between units (turmeric is just always jumpier than wheat --- that
constant gap is soaked up and cannot bias me). Property: constant over time within a unit.
\item $\lambda_t$ --- \textbf{time fixed effects}: one estimated number per month, absorbing
anything common to all commodities that month. This is literally where the Ukraine-war shock
differences out. Property: the same for every unit in a given month.
\item $\varepsilon_{it}$ --- the \textbf{error term}: everything else, unobserved. I do not estimate
it directly; its behaviour governs the standard errors (Section~3.5).
\end{itemize}

\smallskip\noindent\emph{Why the standalone $\mathrm{Treat}_i$ and $\mathrm{Post}_t$ terms vanish.}
They are not written separately because the fixed effects already contain them: $\mathrm{Treat}_i$
never changes over time, so it sits inside $\alpha_i$; $\mathrm{Post}_t$ is the same for everyone in
a month, so it sits inside $\lambda_t$. Only the \emph{interaction} survives --- which is exactly
the term I want.

\applied I ran this twice. On the \emph{national} series (one average price per commodity, 11
commodities, monthly) and on the much richer \emph{district} panel (commodity$\times$district$\times$
month: the assembled panel is about $130{,}000$ rows, and after dropping units too sparse to
estimate a window the headline regression runs on $104{,}491$ observations across $1{,}570$
commodity-district units). The results are in Section~3.6.

\subsection{The assumption that makes it work}

The whole design rests on \textbf{parallel trends}: \emph{absent the ban, treated and control
volatility would have moved in parallel.} I do not need them at the same \emph{level} (the unit
fixed effects $\alpha_i$ handle permanent gaps); I only need the same \emph{trend}. If that holds,
any divergence after the ban is the ban's footprint. If it fails, the DiD is measuring drift, not
the policy. The event study (Section~4) is how I test it.

Three further assumptions, stated plainly: \emph{no anticipation} (behaviour does not jump before
the announcement); that the time fixed effects $\lambda_t$ truly capture the common shock (only
exact if 2022 hit every commodity \emph{equally} --- a weakness I return to); and \emph{no
cross-group spillover} (the jargon is SUTVA --- one unit's treatment must not affect another
unit's outcome). The spillover worry here is concrete: if traders shut out of banned futures
\emph{migrated} into my control crops' futures, the controls would get more volatile precisely
\emph{because} of the ban, contaminating the comparison and biasing the DiD toward zero or negative.
I test that as a separate migration question, and it is one more reason the better-identified design
later picks control weights from the donors \emph{least} entangled with the treated shock.

\subsection{Why I cluster the standard errors by commodity}

\intuit A \emph{standard error} is the margin of error on $\beta$. Volatility within a single
commodity is \emph{autocorrelated} --- today's value is statistically tied to yesterday's; a
turbulent week tends to be followed by another --- so my observations are not independent.
Pretending they are makes the error bars falsely tiny and the result look more certain than it is.
\emph{Clustering by commodity} treats each commodity as one correlated blob and produces honest,
wider error bars.

\smallskip\noindent\emph{The few-clusters caveat.} What governs the reliability of clustered
standard errors is the number of \emph{clusters}, not the number of observations. I have only about
11 commodities $=$ 11 clusters --- and just \emph{five} on the treated side --- which is few.
Crucially this is true \emph{even on the district panel}: $123{,}000$ observations still collapse
into the same $\approx 11$ commodity blobs, so the district result is no better protected than the
national one. With so few clusters even the clustered standard errors are not fully reliable, and
the naive large-sample $p=0.008$ is \emph{over-optimistic} --- I do not treat it as the headline.
The honest fixes are a cluster-robust t-test with the correct (G$-$1) degrees of freedom, which
gives $p\approx0.026$, and a \emph{wild-cluster bootstrap} (a resampling correction that holds up
with few clusters), which gives $p\approx0.046$. Both put the effect right around the
$0.03$--$0.05$ line: \emph{borderline-significant}, not robustly so. And there is a deeper warning
specific to my data: all five treated commodities are \emph{larger} than all five controls (about
an $11\times$ size imbalance), and that is precisely the regime where the wild bootstrap is least
trustworthy (MacKinnon--Webb). So I report the borderline few-cluster numbers as the honest ones
and explicitly flag that even they may flatter the result.

\subsection{What I actually found}

\tells On the \emph{national} series the DiD was essentially zero: $\hat\beta = -0.0115$,
$p = 0.93$ --- no detectable effect, sign unstable, nowhere near the $+8$ to $+10\%$ my very first
rough estimate had suggested. This is a low-power dress rehearsal (only 11 units), so it neither
confirms nor kills anything.

On the much richer \emph{district} panel the current DiD --- with four food-cereal donors in the
control pool --- is $\hat\beta = -0.103$, i.e.\ banned crops' spot volatility about $10\%$
\emph{lower} after the ban relative to controls ($e^{-0.103}-1 \approx -0.098$), and it is
\emph{not} statistically significant under honest few-cluster inference ($p\approx0.145$
cluster-robust t with G$-$1 df; $p\approx0.153$ wild bootstrap). The donor arc matters as much as
the number: at the earlier clean-core stage (industrial/spice controls only) the same regression
read $\hat\beta=-0.232$ ($-20.7\%$) with borderline significance ($p\approx0.026/0.046$; the naive
$p=0.008$ was over-optimistic) --- adding the economically right food controls \emph{halved} the
estimate. This is the headline spec: it drops the MSP-censored paddy (whose government-pinned
price makes its ``volatility'' a mechanical artifact --- see the box below) and replaces the
gum-contaminated guar control (id~75) with the true futures underlying, ``Guar Seed'' (id~413).

\smallskip\noindent\textbf{The check I trust most: leave-one-commodity-out.} An aggregate number
could be one commodity's story in disguise, so I dropped each banned commodity in turn and
re-estimated. The negative sign survives every drop ($-7.4\%$ to $-15.2\%$), and the two extreme
cases are the informative ones. Dropping \emph{wheat} --- the commodity most entangled with the
government support-price (MSP) machinery, the obvious suspect for a spurious result --- makes the
effect \emph{stronger and significant} ($-15.2\%$, $p\approx0.003$): the non-wheat result is not
wheat- or MSP-driven. Dropping \emph{chana} thins the effect to $-7.4\%$ ($p\approx0.29$): the
aggregate signal is substantially a \emph{chana} signal. It is also stable to district-liquidity
and outlier filters ($-8.7\%$ to $-10.0\%$). This decomposition --- not a blanket ``robust'' ---
is what I would actually lead with.

\smallskip\noindent\textbf{Why paddy is dropped (price censoring).} For paddy, $40.3\%$ of daily
spot returns are \emph{exactly zero} --- versus $1$--$6\%$ for the clean commodities --- because
FCI/state procurement pins the price at the Minimum Support Price (MSP) for long stretches. Realized
volatility computed on a government-administered, flat price is a mechanical artifact, not market
volatility: the series is effectively \emph{censored}, so it cannot inform a volatility-effect
estimate and I drop it from the primary spec. (Wheat is also MSP-exposed at $19.7\%$ flat, but it is
a core commodity, so I keep it with an MSP flag and read it as a robustness check, not a clean unit.) A careless reading would shout ``the ban worked, vol fell $21\%$!'' --- but a single DiD
number means nothing until it survives its own stress tests, and the story of \emph{how} this one
came to survive them (Sections~4--5) is the part worth telling. The sign-instability across my
iterations (a first rough $+8$--$10\%$, to roughly zero on the national series, to $-21\%$ on the
clean district panel) is not a mistake to hide; it is exactly why I leaned on falsification and on a
panel of estimators rather than trusting any one before/after.

% =====================================================================
\section{Event study}

\intuit An event study is just DiD \emph{with time resolution}. Instead of one ``after'' dummy that
averages over the whole post-ban period, I split the treatment effect into a \emph{sequence} of
dummies --- one for each month (or bin of months) before and after the ban. Plotting these traces
the banned-vs-control gap month by month, and does two jobs at once. \emph{Job one: it tests
parallel trends.} The pre-ban coefficients (the ``leads'') should all be statistically zero --- if
the groups were truly trending together beforehand, no systematic gap should build up before the
ban. \emph{Job two: it shows the dynamics.} The post-ban coefficients (the ``lags'') trace whether
volatility jumped immediately and stayed high (a real, persistent effect) or merely drifted (noise).

\example Imagine the estimated treated-vs-control gaps, month by month around the ban:
\[
\underbrace{-2:\ +0.01,\quad -1:\ 0\ \text{(baseline)}}_{\text{pre-ban (leads)}},\qquad
\underbrace{+1:\ +0.09,\quad +2:\ +0.02,\quad +3:\ -0.01}_{\text{post-ban (lags)}}.
\]
The pre-ban values sitting near zero would \emph{support} parallel trends; the brief $+0.09$ spike
at month $+1$ that fades is a transitory blip, not a lasting effect. If instead the pre-ban values
were drifting upward (say $-2: +0.05$, building toward the ban), that would \emph{fail} the test ---
the groups were already separating before anything happened.

\formula
\[
\ln\sigma_{it} \;=\; \alpha_i + \lambda_t + \sum_{k\neq -1}\beta_k\,\big(\mathrm{Treat}_i \times \mathbf{1}[t - t^\ast = k]\big) + \varepsilon_{it}.
\]

\vars
\begin{itemize}
\item $\ln\sigma_{it},\ \alpha_i,\ \lambda_t,\ \varepsilon_{it}$ --- exactly as in the DiD
(dependent outcome, unit and time fixed effects, error). Unchanged.
\item $t^\ast$ --- the ban month: a fixed, known date (December 2021), not estimated.
\item $k$ --- ``months relative to the ban'': $k<0$ before, $k>0$ after. An index, not a variable.
\item $\mathbf{1}[t-t^\ast=k]$ --- an \emph{indicator} that is $1$ in exactly the relative month $k$
and $0$ otherwise. Independent variable.
\item $\beta_k$ --- a \emph{whole set} of parameters (one per relative month), each the
treated-vs-control gap in that month. The pre-ban $\beta_k$ ($k<0$) are the parallel-trends test;
the post-ban $\beta_k$ ($k\ge 0$) are the dynamics.
\item $k=-1$ --- the omitted baseline (the month just before the ban). Every other $\beta_k$ is
measured \emph{relative} to it; you must drop one period or the equation is unsolvable (collinear).
\end{itemize}

\applied I estimated the full set of $\beta_k$ on the district panel and ran a \emph{joint test}:
``are all the pre-ban leads zero \emph{together}?'' (a single test across many coefficients at once
--- an $F$-type test). This is a test where I am \emph{rooting for the null}: the null is ``parallel
trends holds,'' which is the assumption I need, so a small $p$-value here is \emph{bad} news.

\tells This test has a history worth telling. On my \emph{first} corrected run the pre-ban leads
were \emph{not} flat --- they drifted, and the joint test \emph{rejected} ($p\approx0.03$), which by
my own pre-registered rules would have retired the DiD. The natural economic story was reverse
causation: SEBI suspended futures precisely because volatility was already high, so banned crops ran
hot into the ban and then mean-reverted, and a naive DiD could read that natural cooling as ``the ban
lowered volatility.'' I took that worry seriously --- but before retiring the design I chased the
\emph{source} of the drift, and it turned out to be a measurement bug, not economics (Section~5). Once
I fixed it, the leads flattened: \textbf{the joint pre-trend test no longer rejects ($p = 0.19$).}
The thing that looked like reverse causation in the leads was largely manufactured by how I had been
measuring volatility on a calendar grid.

\smallskip\noindent\emph{But the joint test masks something, and I won't hide it.} A joint
$p=0.19$ means I cannot reject parallel trends \emph{on average over all the leads} --- it does not
mean every lead is clean. In fact the \emph{most recent} pre-ban bin is still individually
significant: about $+10\%$ ($p=0.049$). So in the months right before the ban the two groups were
already starting to diverge. That residual pre-trend is exactly the kind of thing the
reverse-causation/mean-reversion story predicts, and it is one of the reasons I call the overall
result suggestive rather than conclusive --- the parallel-trends assumption survives the joint test
but is not pristine.

% =====================================================================
\section{The placebo test, and the bug that flipped a dead result alive}

\intuit A placebo test asks: if I run the \emph{identical} analysis at a date where \emph{nothing
happened}, does it still find an ``effect''? If it does, the method is detecting drift, not the
policy. I placed a \emph{fake} ban at 20 December 2019 and used \emph{only} pre-ban data (before the
real December 2021 suspension), so there is no real policy anywhere in the window for it to pick up.

\smallskip\noindent\textbf{The trap to spell out.} This is the one test where I am \emph{rooting for
the null}. The null hypothesis is ``no effect,'' and on a fake date that is the answer I \emph{want}
--- a clean method should detect nothing. So here, unlike a normal test, a \emph{small} $p$-value is
\emph{bad} news: it means my method conjures an effect out of thin air.

\smallskip\noindent\textbf{The bug --- my favourite part of the project.} On my \emph{first}
corrected run this placebo \emph{failed}: it returned a ``significant'' $-13.6\%$ ($p=0.046$) at a
fake date, and the joint pre-trend test rejected too --- two independent checks both screaming that
the design was detecting drift, not the policy. By my pre-registered rules the DiD was dead. Rather
than bury it, I went looking for \emph{why} a fake date found an effect, and traced it to a
\emph{calendar-grid measurement bug}. I had laid the mandi spot series on a full calendar (every
day, weekends and holidays included), and on non-trading days the price simply carried forward
unchanged. But my realized volatility annualised on $252$ \emph{trading} days. So the carried-forward
weekend prices injected runs of artificial zero-returns and stale levels, manufacturing
\emph{spurious autocorrelation} in the volatility series --- exactly the kind of slow, persistent
drift that makes a placebo light up and pre-trends bend. The fix was unglamorous: filter to weekdays
(actual trading days) before computing returns, so the $252$-day annualisation matches the grid the
data actually live on. I also dropped the MSP-censored paddy at the same time. Together those two
changes cleaned up the falsification battery.

\tells After the fix, the placebo is \emph{no longer significant} --- and it improved in two
stages worth narrating. At the clean-core stage (industrial/spice donors) it was a \emph{marginal}
pass: the fake Dec-2019 ban returned $-11.7\%$ ($p = 0.12$ analytic, $0.19$ bootstrap), about 56\%
of the then-headline $-20.7\%$ --- a warning that part of the ``effect'' could be drift or
mean-reversion after the pre-ban run-up, with a hot final pre-bin ($+10\%$, $p=0.049$) lurking
under the joint pre-trend test. On the \emph{current food-donor panel} the battery is genuinely
clean: the placebo returns $\hat\beta = -0.067$ ($-6.5\%$) with $p = 0.34$ analytic, $0.43$ on the
bootstrap, and the joint pre-trend test is comfortable ($p = 0.45$) with \emph{no} individually
significant lead bin left. Notice what happened: with the economically right counterfactual, the
placebo and the headline shrank \emph{together} ($-6.5\%$ vs.\ $-9.8\%$) --- the drift the old
placebo was picking up lived in the bad donor pool, not the design. This is still the
research-integrity story I would tell an interviewer --- my first corrected DiD looked dead, I
refused to paper over it, I found a measurement bug, then let a better donor pool discipline the
estimate --- but I am explicit that the surviving aggregate result is \emph{suggestive, not
conclusive}: it is small ($-9.8\%$) and not significant once the proper food-staple donors are in.

% =====================================================================
\section{Synthetic control and Synthetic DiD: corroborating the point estimate}

A single DiD point estimate is more believable if other estimators land near it, so alongside the
DiD I built a panel of comparative-case methods. They do all come out negative, in the same
$-20$ to $-27\%$ range. But I want to be careful here, because it is tempting --- and wrong --- to
sell this as ``four independent methods converge.'' For reasons I lay out at the end of the section,
the agreement is partly \emph{mechanical} (same panel, same five-and-five commodities), the
synthetic-DiD pooled number has \emph{no valid standard error}, and the classic synthetic controls
fit poorly. So the honest framing is: these corroborate the \emph{sign and rough magnitude} of the
DiD point estimate; they are not four independent confirmations of significance.

\subsection{The idea: build the missing twin}

\intuit DiD uses an \emph{equally-weighted} average of the control crops as the counterfactual ---
``what banned crops would have done absent the ban.'' But maybe no single control is a good stand-in
for chana, while a \emph{blend} --- say $0.5$ guar $+ 0.3$ jeera $+ 0.2$ castor --- tracks chana's
pre-ban path almost exactly. \emph{Synthetic control} (Abadie) finds exactly that blend: it picks
weights on the donor crops so the weighted donor matches the treated crop's \emph{pre-ban}
volatility as closely as possible, then carries that ``synthetic twin'' forward past the ban. The
gap between the real treated series and its synthetic twin after the ban is the estimated effect.
\emph{Synthetic DiD} (Arkhangelsky) is the modern hybrid: it keeps DiD's two-way structure but lets
\emph{both} the unit weights (which donors) \emph{and} the time weights (which pre-periods) be chosen
by the data, so it is more robust than either plain DiD or plain synthetic control alone.

\formula In words rather than full notation (the mechanics are weighted least squares with a
pre-fit constraint): synthetic control chooses donor weights $w_j \ge 0$, $\sum_j w_j = 1$, to
minimise the pre-ban distance between the treated path and $\sum_j w_j \times (\text{donor }j)$;
the post-ban treated-minus-synthetic gap is the effect. Synthetic DiD augments the DiD regression
with those data-chosen unit weights $\hat w_i$ and time weights $\hat\lambda_t$:
\[
(\hat\beta^{\text{sdid}}) \;=\; \arg\min \sum_{i,t} \hat w_i\,\hat\lambda_t\,\big(\ln\sigma_{it} - \alpha_i - \lambda_t - \beta\,\mathrm{Treat}_i\!\times\!\mathrm{Post}_t\big)^2.
\]

\applied I ran all four and report the panel, with the inference stated honestly for each:
\begin{itemize}
\item \textbf{Two-way-FE DiD} (the headline, Section~3): $-9.8\%$ on the current food-donor panel;
\emph{not significant} under few-cluster inference ($p\approx0.145$ cluster-robust t,
$\approx0.153$ wild bootstrap). Ex-wheat it is $-15.2\%$ ($p\approx0.003$); ex-chana it thins to
$-7.4\%$ --- the signal is concentrated in chana.
\item \textbf{Synthetic DiD} (pooled): $-19.1\%$, and with nine donors the placebo SE is now
\emph{valid}: $z=-1.88$, borderline. (An earlier ``$z=-7.7$'' I had quoted was a donor-jackknife
\emph{artifact} and I have removed it; at the five-donor stage the pooled SE was undefined
altogether.) Per commodity, chana is the one significant case ($-38.7\%$, $z=-2.05$); wheat is
$\approx$zero ($+0.4\%$); mustard, soybean and moong are negative but insignificant.
\item \textbf{Commodity-level synthetic control} (one synthetic twin per banned crop, on
district-median paths): chana $-37.3\%$ \emph{at} the placebo floor ($p=0.100$ with nine donors);
mustard $-28.2$, wheat $-37.8$, soybean $-27.6$ carry weak placebo p's ($0.3$--$0.6$), and moong is
small ($-8.7\%$). Corroborates chana; weak elsewhere.
\item \textbf{District-unit synthetic control} (one synthetic per treated district): all treated
commodities come out negative, strongest for chana/mustard/wheat (each at the $p=0.100$ floor).
The floor is $1/10$ with nine donors --- better than the old $1/6$, but still \emph{structurally
unable} to reject at $5\%$; the ``significance'' remains capped by construction.
\end{itemize}

\tells \textbf{The estimators agree on sign --- banned-commodity spot volatility came out
\emph{lower}, not higher, the opposite of the going $+8$--$10\%$ hypothesis --- but the aggregate
effect is modest ($\approx-10\%$ DiD, $-19\%$ SDID) and not robustly significant; the strong,
consistent signal is chana.} I will not oversell the agreement, and three honest caveats are why.
\emph{First}, it is partly \emph{mechanical}: DiD, SDID and SCM all use the same panel, outcome and
treated set, and the SDID unit weights come out roughly uniform, so SDID collapses back toward the
plain DiD --- one comparison re-weighted, not four independent confirmations. \emph{Second}, the
inference is genuinely weak outside chana: the pooled SDID $z=-1.88$ is borderline and the SCM
placebo p's are mostly $0.3$--$0.6$. \emph{Third}, the donor arc is the deepest lesson: swapping the
industrial/spice pool for food-cereal donors \emph{halved} the headline ($-20.7\% \to -9.8\%$) and
removed its significance --- the estimate answers to the quality of its counterfactual, and the
remaining oilseed donors (groundnut, sesamum, sunflower) are still to land, so mustard/soybean are
provisional. So I frame this as: \emph{the ban did not raise volatility; the aggregate decline is
suggestive, not conclusive; the defensible positive evidence is chana}.

% =====================================================================
\section{GARCH(1,1)}

The DiD compares \emph{across} commodities. GARCH instead looks \emph{inside} one commodity's own
price process over time and asks whether its volatility regime changed at the ban. It is the
standard tool for modelling how volatility evolves --- and a different kind of test: \emph{no control
group at all}, just one commodity against its own history.

\subsection{Volatility clustering}

\intuit One robust fact about prices --- stocks, currencies, or chana in a mandi --- is that calm
and turbulent days come in \emph{clumps}: a wild day tends to be followed by more wild days, a
sleepy one by more sleepy ones. Markets have ``stormy seasons'' and ``quiet seasons'' that persist
before flipping. This is \emph{volatility clustering}. A model that assumes one fixed variance for
the whole sample --- as if every day had identical weather --- misses it entirely. GARCH
(Generalised Autoregressive Conditional Heteroskedasticity) lets today's variance depend on what
just happened, so it captures the clustering. ``Conditional'' means today's variance depends on
recent history; ``heteroskedasticity'' means the variance changes over time instead of staying
constant.

\subsection{The model}

\formula GARCH(1,1) describes the return $r_t$ and its time-varying variance $\sigma_t^2$:
\[
r_t = \mu + \varepsilon_t,\qquad \varepsilon_t = \sigma_t z_t,\quad z_t \sim \text{iid}(0,1),
\]
\[
\sigma_t^2 \;=\; \omega \;+\; \alpha\,\varepsilon_{t-1}^2 \;+\; \beta\,\sigma_{t-1}^2.
\]
In words: \emph{today's variance $=$ a baseline constant $+$ a slice of yesterday's squared
surprise $+$ a slice of yesterday's variance.}

\vars
\begin{itemize}
\item $r_t$ --- the daily return. \emph{Observed data}, the input series (here, one commodity's
daily spot returns). In the first line it is the dependent variable of the ``mean equation.''
\item $\mu$ --- the average return. A \emph{parameter} (estimated), usually tiny.
\item $\varepsilon_t$ --- the ``surprise'' or shock: the part of the return not explained by the
mean ($\varepsilon_t = r_t - \mu$). A \emph{derived, unobserved} quantity.
\item $z_t$ --- a standardized random innovation, ``iid$(0,1)$'' meaning independent draws with mean
$0$ and variance $1$. The pure randomness; not estimated.
\item $\sigma_t^2$ --- the \emph{conditional variance} (today's variance given the past). This is a
\emph{latent} variable: not directly observed, produced \emph{by} the model. It is the real object
of interest --- volatility is $\sigma_t = \sqrt{\sigma_t^2}$.
\item $\omega$ (omega) --- a \emph{parameter}: the small constant \emph{intercept} of the variance
equation. Property/trap: it is \emph{not} the long-run volatility level itself (a common slip ---
see below). Constraint: $\omega>0$ so variance stays positive.
\item $\alpha$ (alpha) --- a \emph{parameter}, the \emph{ARCH term}: the weight on yesterday's
squared surprise $\varepsilon_{t-1}^2$. High $\alpha$ $=$ a jumpy, news-sensitive series that reacts
hard to a shock. Constraint: $\alpha \ge 0$.
\item $\beta$ (beta) --- a \emph{parameter}, the \emph{GARCH term}: the weight on yesterday's
\emph{variance} $\sigma_{t-1}^2$. High $\beta$ $=$ long memory; once stormy, it stays stormy.
Constraint: $\beta \ge 0$. (Note: this $\beta$ is unrelated to the DiD's $\beta$ --- same Greek
letter, different model.)
\end{itemize}

\subsection{Persistence and the long-run level}

\intuit The key quantity is the sum $\alpha + \beta$, called \emph{persistence}: when a shock hits,
how long does the elevated volatility last? Near $1$ it decays very slowly (a storm today still
raises volatility weeks out); lower, and shocks fade fast. The model is only \emph{stationary} ---
meaning the variance settles to a stable long-run level instead of exploding --- when
$\alpha + \beta < 1$. That inequality is the central \emph{property} of the parameters. The
\emph{long-run (unconditional) variance} the series reverts to is
\[
\sigma^2_{\text{long-run}} \;=\; \frac{\omega}{1 - \alpha - \beta},
\]
which depends on $\omega$ \emph{and} the persistence together --- so $\omega$ alone is not the level.

\example \emph{Half-life intuition.} ``Half-life'' $=$ how many days for a volatility shock to lose
half its extra height; the formula is $\ln(0.5)/\ln(\text{persistence})$. If persistence $=0.99$,
half-life $\approx \ln(0.5)/\ln(0.99) \approx 69$ days --- about ten weeks; storms linger. If
persistence $=0.78$, half-life $\approx \ln(0.5)/\ln(0.78) \approx 2.8$ days --- shocks die in days
and the market settles quickly. This is the arithmetic that \emph{would} make a drop in persistence
economically meaningful --- a market going from weeks-long storms to days-long ones. I flag it here
because an early, dirty fit of my data showed exactly such a drop ($0.99 \to 0.78$) and looked like a
striking ``market calmed'' result; on the trading-day-clean data it does \emph{not} reproduce (the
clean persistences are roughly $0.55 \to 0.69$, going the other way, with volatility about flat), so
I do not lean on it. The arithmetic is sound; the dramatic input numbers were an artifact.

\subsection{Testing the ban: a step dummy}

\formula To ask directly ``did the ban shift the volatility \emph{level}?'', I add a post-ban
\emph{step dummy} to the variance equation:
\[
\sigma_t^2 \;=\; \omega + \alpha\,\varepsilon_{t-1}^2 + \beta\,\sigma_{t-1}^2 + \gamma\,\mathrm{Post}_t.
\]

\vars
\begin{itemize}
\item $\mathrm{Post}_t$ --- an \emph{independent} indicator: $0$ before 20 December 2021, $1$ after.
A known, fixed switch, not estimated.
\item $\gamma$ (gamma) --- the \emph{parameter} of interest here: how much the baseline variance
\emph{steps} at the ban. A positive, significant $\gamma$ $=$ volatility stepped \emph{up} at the
ban (over and above normal clustering); negative $=$ stepped \emph{down}.
\item All other terms as above; $\alpha,\beta$ still drive the day-to-day clustering \emph{around}
the shifted baseline. (In estimation $\mathrm{Post}_t$ enters as an exogenous regressor and I keep
the implied variance constrained positive.)
\end{itemize}

\applied I fit GARCH(1,1) per commodity and split the sample at the ban, and I let the \emph{data}
choose the break date with an ICSS (Iterated Cumulative Sums of Squares) variance-break detector ---
the honest version of the test, since imposing the ban date is exactly the thing in question. The
result on clean data is a \emph{null}: the ICSS detector finds \emph{no} variance break near the
suspension on any series. The cleanest banned staple, \emph{chana}, has roughly flat unconditional
volatility across the ban (its persistence even nudges \emph{up}, $\approx 0.55 \to 0.69$), and the
district-median chana fit is close to degenerate (persistence pinned near $1$), so I do not read fine
movements into it. An early, calendar-dirty fit had suggested a dramatic calming (persistence
$0.99 \to 0.78$, vol $71\% \to 46\%$); that was the same measurement artifact the DiD placebo exposed,
and it washes out on trading-day-clean returns.

\tells \textbf{C2-volatility is inconclusive, and the volatility verdict rests on C1.} The original
``volatility rose, significant at $10\%$'' claim does \emph{not} reproduce --- but neither does a
clean ``the market calmed'' counter-result hold up inside the GARCH: with no data-chosen break and
roughly flat conditional volatility, the within-commodity time-series evidence is simply silent. That
is the honest read, and it is why I lean the whole volatility case on the cross-commodity panel
(chiefly the leave-one-out-robust DiD point estimate, corroborated in sign by Synthetic DiD and
synthetic control), rather than on GARCH.
For completeness, pinning any regime change to the policy would still want a model that lets the data
choose the break --- a Bai--Perron endogenous structural-break test or a Markov-switching GARCH ---
but on this data there is no break for it to find.

% =====================================================================
\section{The basis}

The last concept ties futures and spot together --- and is where one part of the original project
framing turns out to be internally impossible for the banned crops.

\subsection{Cost of carry and convergence}

\intuit The \emph{basis} is the difference between the futures price and the spot price of the same
commodity at the same moment. Why are they ever different? \emph{Cost of carry.} If you want chana
in three months, you can buy a futures contract, or you can buy today and store it --- and storing
costs warehousing, insurance, and forgone interest. So the futures price for later delivery normally
sits \emph{above} today's cash price by roughly those carrying costs. That gap is the basis, and a
positive basis (futures above spot) is called \emph{contango}. Agricultural commodities, though,
often flip to \emph{backwardation} (futures \emph{below} spot) in a near-term supply squeeze, when
having the physical crop right now is worth a premium --- that premium is the \emph{convenience
yield} (the value of holding the actual crop rather than a paper claim; e.g.\ a mill that cannot
afford to run out). As a futures contract nears delivery, the ``three months of storage'' shrinks
toward ``zero,'' so the carrying-cost gap shrinks and futures and spot \emph{converge}: the basis
goes to about zero at expiry. Arbitrage enforces this --- if they did not converge, a trader could
lock in a near-riskless profit, and that very trading pulls them together.

\formula The basis itself is simply
\[
\text{basis}_t \;=\; F_t - S_t,
\]
and the fuller \emph{cost-of-carry} relation that explains \emph{why} $F_t$ and $S_t$ differ is
\[
F_t \;=\; S_t \, e^{(r + u - y)\,T}.
\]

\vars
\begin{itemize}
\item $S_t$ --- the spot (cash) price now. \emph{Observed data}.
\item $F_t$ --- the futures price now, for delivery in $T$ years. \emph{Observed data}.
\item $\text{basis}_t$ --- their difference, $F_t - S_t$. \emph{Derived}. (Some practitioners define
it spot minus futures --- a pure sign convention; I use futures minus spot.) A \emph{widening} basis
$=$ futures and spot drifting apart by more than carry justifies, usually read as market stress or
impaired arbitrage --- the price-discovery link breaking down.
\item $r$ --- the financing (interest) rate, an annual rate. Pushes $F$ \emph{above} $S$.
\item $u$ --- the storage cost, an annual rate. Also pushes $F$ \emph{above} $S$.
\item $y$ --- the convenience yield, an annual rate. Pulls $F$ \emph{below} $S$. The term with no
equity analogue --- it is why \emph{agricultural} futures behave so differently from the index
futures a trading desk usually sees. When $y$ is large (a tight physical market), the exponent can go
negative and futures fall below spot: that is backwardation.
\item $T$ --- time to delivery, in years. As $T \to 0$ the whole exponent $\to 0$, so $F_t \to S_t$:
that is convergence at expiry.
\end{itemize}
(Equity-index futures are the clean special case where $u=0$ and $y$ is just the dividend yield, so
they almost never sit in deep backwardation.)

\example Spot chana is $\rupee\,5{,}000$ per quintal today; the three-month futures trades at
$\rupee\,5{,}150$. The basis is $F_t - S_t = 5{,}150 - 5{,}000 = +150$, i.e.\ $+3\%$ --- contango,
reflecting about three months of storage and financing. As delivery approaches, that $+150$ should
shrink toward $0$. If instead the futures were at $\rupee\,4{,}900$, the basis would be $-100$:
backwardation, signalling a near-term scarcity (a high convenience yield).

\subsection{A note on rollover}

\intuit A futures contract \emph{expires}, so to track a continuous futures price past expiry you
must \emph{roll}: close the contract about to expire and open the next one out. I keep the
individual nearby legs (the first-, second-, and third-to-expire contracts, c1/c2/c3) precisely so I
control this. It matters because rolling is not free: in contango you sell the cheap expiring
contract and buy the dearer next one and \emph{bleed} that gap (negative roll yield); in
backwardation you earn it. For my \emph{spot}-volatility outcome rollover does not contaminate
anything --- the spot mandi series has no contracts to roll. But it would distort any futures-return
or basis series built from a single naively spliced price, which is why I keep the separate
c1/c2/c3 legs rather than one blended series.

\subsection{Why a banned commodity has no post-ban basis}

\tells Here is the clean logic point. The basis needs \emph{both} a futures price and a spot price.
The suspension switched off \emph{futures} for exactly the banned commodities after 20 December 2021.
After that date there is simply \emph{no futures price to subtract}, so a ``post-ban basis'' for a
banned commodity is \emph{undefined} --- it cannot exist. Any number reported under that label must
be measuring something else (most likely stale, zero-volume settlement marks that some contracts
kept publishing for months after trading stopped --- a trap I guard against by truncating banned
futures at the ban date). What I \emph{can} honestly measure is the basis only where it is defined:
the \emph{pre-ban} basis trend for the banned crops (futures were live then), and the
\emph{full-window} basis for the control crops (which kept trading).

\applied The basis even served as a \emph{data-quality instrument}. My first guar spot series
correlated only $0.04$ with guar futures and produced a nonsensical basis --- which is how I caught
that I had grabbed a gum-contaminated series (guar \emph{gum}, a processed product about ten times
the price of guar \emph{seed}). Switching to the true ``Guar Seed'' underlying restored a $0.99$
correlation and a sensible, stable basis. A quantity that should track its own futures, but does
not, is a red flag worth chasing.

% =====================================================================
\section{How the pieces fit together}

Each method answers a different slice of the same question, and reading them together is what makes
the project honest:
\begin{itemize}
\item \textbf{Realized volatility} turns raw prices into a clean, model-free ``jumpiness'' number ---
the outcome everything else explains. Getting its calendar grid right (trading days, not weekends
carried forward) turned out to be load-bearing for the entire result.
\item \textbf{DiD} compares the change in that number for banned vs.\ still-trading crops, so common
2022 shocks difference out. The current estimate, with four food-cereal donors in the pool, is
$-9.8\%$ --- \emph{not significant} under few-cluster inference ($p\approx0.145$ cluster-robust t,
$\approx0.153$ wild bootstrap). The donor arc is the real lesson: on the earlier industrial/spice
pool the same DiD read $-20.7\%$ with borderline significance; the economically right counterfactual
\emph{halved} it. Leave-one-out shows the signal is \emph{concentrated in chana}: ex-wheat is
$-15.2\%$ ($p\approx0.003$, significant), while ex-chana falls to $-7.4\%$.
\item \textbf{The event study and placebo} stress-test that estimate. On my first corrected run they
\emph{failed} (pre-trend $p\approx0.03$, a fake 2019 ban ``significant'' at $-13.6\%$) --- but chasing
that failure exposed a calendar-grid measurement bug, not reverse causation. With the bug fixed the
pass was marginal (placebo $-11.7\%$, $\approx$56\% of the then-headline; a hot final pre-bin); on
the food-donor panel it is genuinely clean --- placebo $-6.5\%$ ($p\approx0.34/0.43$), joint
pre-trend $p\approx0.45$, no significant lead bin. The design survives the tests that nearly
retired it, and the tests got cleaner as the counterfactual got better.
\item \textbf{Synthetic control and Synthetic DiD} build a data-weighted ``twin'' for each banned
crop. Synthetic DiD pools to $-19.1\%$ and, with nine donors, now has a \emph{valid} placebo SE
($z=-1.88$, borderline); per commodity, \textbf{chana is the one significant case ($-38.7\%$,
$z=-2.05$)}, wheat is $\approx$zero ($+0.4\%$ --- its earlier decline was the MSP/export-ban
confound), and the rest are negative but insignificant. Commodity SCM: chana $-37.3\%$ at the
$1/10$ placebo floor; others weak. I do \emph{not} read this as ``four independent estimators
converge'': they share the same panel and treated set, so the agreement is partly mechanical. It
supports the sign; it does not certify aggregate significance.
\item \textbf{GARCH(1,1)} looks inside a single commodity over time --- no control group. It neither
revives the original ``volatility rose'' claim (an ICSS detector finds \emph{no} variance break at
the ban) nor cleanly supports the opposite (conditional volatility is roughly flat once the calendar
bug is removed). C2-volatility is therefore inconclusive, and the volatility case rests on the C1
panel above.
\item \textbf{The basis} shows why one part of the original framing is internally impossible for
banned crops, and doubles as a data-quality check.
\end{itemize}
The through-line: the going hypothesis was that switching off futures would \emph{raise} spot
volatility ($+8$--$10\%$); my analysis \emph{refutes} that --- the point estimate is firmly
negative, around $-20\%$, and robust to dropping any single commodity, including the MSP-flagged
wheat. But the calibrated claim is deliberately hedged: \emph{suggestive, not conclusive.} The
significance is borderline under honest few-cluster inference, the placebo recovers about half the
effect, the estimator agreement is partly mechanical rather than four independent confirmations, the
SCM fits poorly, the donor pool is weak, and the ban's reactive timing leaves a mean-reversion threat
the marginal placebo cannot fully rule out. The most interview-worthy moment is not the headline but
how I earned it --- my first corrected DiD looked dead, I refused to paper over a failing placebo, I
traced it to a calendar-grid measurement bug, and fixing it flipped a ``dead'' result into a live (if
borderline) one. A project willing to let its own pre-registered tests nearly kill its headline ---
and to state plainly where the evidence is still merely suggestive --- is doing what empirical work
is supposed to do.

\end{document}
